\scanpage{168}
\textcircled{b} $\varphi_0(T)\in S'^+$ \deleted{\unsure{quelc.}} mais sera fixé par condition 2)
\[
 \deleted{\varphi_0(T)=\sum_{i\geq1}\deleted{\unsure{\alpha_i}}\gamma'_iT^{(i)}\qquad(\gamma'_i\in k',\ i\geq1)}
\]
\deleted{\textcircled{2}} NB si on suppose qu'on se donne déjà $\psi=\varphi_{00}:k\longrightarrow k'$, et qu'on exige $\varphi_*$ $\psi$-semilinéaire i.e. $\varphi_0$ $\psi$-semi-linéaire, alors la donnée 1) se réduit à (b).

\textcircled{2}
\[
 A'_1\in S'_{\deleted{\unsure{*}}}\{X_0\}\simeq k'\{X_0,X_1\}
 =\sum_{\substack{\unsure{i,j\geq0}\\\deleted{\unsure{i\geq0}}\\j\geq0}}\alpha'_{ij}X_0^{(i)}X_1^{(j)}
\]
La condition $A'_1\in S'\{X_0\}^{++}$ signifie $\alpha'_{ij}=0$ si $i=0$
i.e.
\[
 A'_1=\sum_{\substack{i\geq1\\j\geq0}}\alpha_{ij}X_0^{(i)}X_1^{(j)}.
\]
On a
\[
 \beta_{S'}(A'_1)=\sum_{i\geq1}\alpha_{i,0}T^{(\unsure{j|i})},
\]
et on aura donc
\[
 \varphi_0(T)=\sum_{i\geq1}\alpha_{i,0}T^{(i)}
\]
\textcircled{3} $B'_i\in S'\{X_0\}\simeq k'\{X_0,X_1\}$ sont liés par la seule condition
\[
 (2B'_i)J=0.
\]
En résumé :

\uline{Cor. 5} Les hom. d'anneaux diff $\tfrac12$-simpliciaux
\[
 \mathcal A_{k*}=\mathcal A(k\{T\},k\{T\}^+,T)\mathbin{\text{---}}\mathcal A_{k'*}=\mathcal A(k'\{T\},k'\{T\}^+,T)
\]
sont définis par les systèmes \unsure{d'objts}
\[
 \left\{\begin{aligned}
 &\varphi_{00}:k\mathbin{\text{---}}k'\{T\}\quad\text{hom d'anneaux}\\
 &A'_{\deleted{1}}\in\deleted{\unsure{S'}}k'\{T\}\{X_0\}^{\deleted{++}}\simeq k'\{X_0,X_1\}\\
 &B'_i\in k'\{T\}\{X_0\}\simeq k'\{X_0,X_1\}\qquad(i\geq1)
 \end{aligned}\right.
\]
satisfaisant les conditions
\[
 \begin{aligned}
 a)\quad&A'_{\deleted{1}}\in k'\{T\}\{X_0\}^{++}\quad\text{i.e. }A'_{\deleted{1}}=\sum_{\substack{i\geq1\\j\geq0}}\alpha'_{ij}X_0^{(i)}X_1^{(j)}\\
 b)\quad&(2B'_i)J'=0\quad(\forall i\geq1)
 \end{aligned}
\]
où $J'\subset S_{k'}=k'\{T\}$ est l'idéal de $S_{k'}$ engendré par les coeff. $a'_j$ et $b'_{ij}$
\deleted{\unsure{coefficients de $A'$ et des $B'_i$}}
de
\[
 A'=\sum_{j\geq1}a'_jX_0^{(j)}\quad(a'_j\in k'\{T\}=S_{k'})
\]
et des
\[
 B'_i=\sum_{j\geq0}b'_{ij}X_0^{(j)}\quad(b'_{ij}\in k'\{T\}=S_{k'}).
\]
Alors b) équivaut aussi à
\[
 (b')\quad(2b'_{ij})J'=0\quad\text{pour }i\geq1,\ j\geq0.
\]
\margin{167}
